\documentclass[11pt]{article}
\usepackage[a4paper,margin=1in]{geometry}
\usepackage{amsmath,amssymb,amsthm,mathtools}
\usepackage{graphicx}
\usepackage{booktabs}
\usepackage{hyperref}
\usepackage{enumitem}
\usepackage{microtype}
\usepackage{xcolor}
\graphicspath{{../software/origin_iii/figures/}}
\hypersetup{colorlinks=true,linkcolor=blue,citecolor=blue,urlcolor=blue}

\newtheorem{definition}{Definition}
\newtheorem{proposition}{Proposition}
\newtheorem{theorem}{Theorem}
\newtheorem{lemma}{Lemma}
\newtheorem{remark}{Remark}
\newtheorem{example}{Example}

\newcommand{\Ztwo}{\mathbb Z_2}
\newcommand{\sgn}{\operatorname{sgn}}
\newcommand{\spec}{\operatorname{spec}}
\newcommand{\diag}{\operatorname{diag}}
\newcommand{\lambdaMin}{\lambda_{\min}}

\title{Boundary Compensation Origin III:\\
Non-Factorizable Holonomy and Frustrated Multi-Shadow Spectral Geometry\\
\large Edge Signs, Cycle Products, Gauge Classes and Horizon Diagnostics in Structural Shadow Operators}
\author{A. A. Malachevsky\\
Independent researcher\\
Boundary Compensation Origin Program (BC-Origin / BC-O)}
\date{Version v0.1.0-manuscript draft, June 2026}

\begin{document}
\maketitle

\begin{abstract}
BC-Origin I introduced a minimal constructive toy model in which observable shadows arise from hidden oriented residual generators. In that model a hidden index-winding pair $h=(q,n)$ produces a scale channel $|n|$, an orientation channel $\sgn(n)$, a structural scale denominator, and a signed two-shadow spectral operator. BC-Origin II replaced the free pairwise coupling constant $\kappa$ by a structural kernel $K_{ij}=K(n_i,n_j)$ on hidden winding-index space. It also identified a limitation of the baseline sign rule: the signs $s_{ij}=\sgn(n_i n_j)$ are vertex-factorizable and therefore balanced on every cycle. Consequently, BC-Origin II does not produce genuine cycle-level frustration without additional non-factorizable edge data.

The present paper develops the next finite-dimensional toy-model layer. We introduce a symmetric edge-holonomy field $\varepsilon_{ij}\in\{-1,+1\}$ on the coupling graph and distinguish vertex-factorizable signs from non-factorizable cycle data. Under vertex reorientation, $\varepsilon_{ij}\mapsto g_i\varepsilon_{ij}g_j$, while the cycle product
\[
\Omega_C=\prod_{(ij)\in C}\varepsilon_{ij}
\]
is invariant. A sign field is balanced precisely when all cycle products are $+1$; a cycle with $\Omega_C=-1$ carries a non-removable $\mathbb Z_2$ holonomy and is frustrated in the finite-dimensional graph sense used here.

The baseline BC-Origin III operator is the holonomy-weighted off-diagonal multi-shadow operator
\[
(D^{\rm hol}_N)_{ii}=d_i,\qquad
(D^{\rm hol}_N)_{ij}=\eta\,\chi_{ij}K_{ij},\quad i\ne j,
\]
where $\chi_{ij}=\sigma_i\varepsilon_{ij}\sigma_j$ and $\sigma_i=\sgn(n_i)$. This convention keeps the core spectral statements gauge-clean: vertex reorientation changes the matrix by orthogonal conjugation and hence preserves its spectrum. For an identical three-shadow complete graph with edge magnitude $k$, the balanced triangle has spectral type
\[
d+2\eta k,\quad d-\eta k,\quad d-\eta k,
\]
whereas a frustrated triangle has
\[
d+\eta k,
\quad d+\eta k,
\quad d-2\eta k.
\]
Thus non-factorizable holonomy changes the spectral class of the multi-shadow operator and changes the admissibility-horizon margins of the toy shadows.

The construction is a finite-dimensional structural model. It is not a physical lattice gauge theory, not a spin-glass model of matter, not a derivation of particle physics, not a theory of physical time, and not an empirical claim. Its controlled contribution is to show how the structural-kernel shadow geometry of BC-Origin II acquires genuine cycle-level obstruction once non-factorizable edge holonomy is admitted.
\end{abstract}

\paragraph{Keywords.}
Boundary Compensation; BC-Origin; structural kernels; edge holonomy; graph frustration; finite-dimensional spectral geometry; hidden winding index; admissibility horizon; toy model; gauge equivalence.

\paragraph{Software repository.}
The accompanying computational companion is intended for the public repository
\url{https://github.com/AIDevelopersMonster/boundary-compensation}.

\paragraph{Author responsibility note.}
The author is the sole responsible scientific decision maker. AI-assisted tools may have been used for technical drafting, structuring, consistency checks, software scaffolding, and LaTeX preparation. AI-assisted tools are not authors or coauthors.

\tableofcontents

\section{Introduction}

The BC-Origin branch asks whether observable elementary structures can be modeled as non-fundamental shadows of hidden residual generators. The first two modules established the minimal static architecture.

BC-Origin I introduced the hidden generator
\begin{equation}
 h=(q,n),\qquad n\in\mathbb Z\setminus\{0\},
\end{equation}
with two observable channels:
\begin{equation}
 n\longmapsto (|n|,\sgn(n)).
\end{equation}
The scale channel enters a closure denominator
\begin{equation}
 d_i=2\pi |n_i|-\theta(|q_i|),
\end{equation}
and the orientation channel enters a signed two-shadow operator. In that minimal model, observable scale ratios are interpreted as
\begin{equation}
 S_a=\frac{\mu_{\rm eff}}{\lambda_a},\qquad \lambda_a>0,
\end{equation}
where $\lambda_a$ is an eigen-denominator. The admissibility horizon is the condition $\lambda_a=0$: as $\lambda_a\to0^+$ the scale ratio diverges and the branch no longer defines a localized shadow inside the toy model.

BC-Origin II addressed the main weakness of the first model. The off-diagonal coupling $\kappa$ should not be chosen independently for every pair of shadows. The central replacement was
\begin{equation}
 \kappa\longrightarrow \kappa_{ij}=K(n_i,n_j),
\end{equation}
where $K$ is a declared structural kernel on hidden winding-index space. This replacement converts pairwise numerical fitting into a structural rule.

However, BC-Origin II also revealed a new limitation. Its baseline sign rule
\begin{equation}
 s_{ij}=\sgn(n_i n_j)=\sigma_i\sigma_j,
 \qquad \sigma_i:=\sgn(n_i),
\end{equation}
is vertex-factorizable. Every cycle product is $+1$. Therefore the baseline graph is balanced and cannot produce genuine cycle-level frustration. The present paper closes exactly this gap.

The new object is an edge-holonomy field
\begin{equation}
 \varepsilon_{ij}=\varepsilon_{ji}\in\{-1,+1\},
\end{equation}
which is not assumed to be expressible as a product of vertex signs. This is the minimal additional structure needed to make the multi-shadow graph carry a non-removable cycle obstruction.

\section{Scope and non-scope}

The paper is deliberately narrow. It studies finite-dimensional real symmetric toy operators on a finite graph of shadows. Words such as ``holonomy'', ``gauge'', and ``frustration'' are used in their graph-theoretic and finite-dimensional toy-model sense.

The paper does not claim:
\begin{itemize}[leftmargin=2em]
\item to construct physical gauge fields;
\item to derive a real spin-glass material;
\item to produce particle masses or Standard-Model data;
\item to define physical time or dynamics;
\item to validate the BC-Origin program empirically.
\end{itemize}

The positive claim is smaller and sharper: non-factorizable edge holonomy is the minimal graph-level extension of BC-Origin II required for genuine cycle obstruction in the structural multi-shadow operator.

\section{Hidden winding labels, kernels and baseline signs}

Let
\begin{equation}
 I_N=\{n_1,\ldots,n_N\}\subset \mathbb Z\setminus\{0\}
\end{equation}
be a finite set of hidden winding labels. A declared kernel gives a symmetric matrix
\begin{equation}
 K_{ij}=K(n_i,n_j),\qquad K_{ij}=K_{ji},\qquad K_{ii}=1.
\end{equation}
Examples include the index-space Lorentzian kernel
\begin{equation}
 K_{ij}=\frac{1}{1+(n_i-n_j)^2}
\end{equation}
and the exponential kernel
\begin{equation}
 K_{ij}=e^{-\alpha |n_i-n_j|}.
\end{equation}
These are declared index-space kernels unless an independent hidden geometry is supplied. In particular, the Lorentzian expression should not be confused with an off-diagonal Fourier matrix element of $(1-\Delta_{S^1})^{-1}$ between pure Fourier modes.

The baseline orientation signs are
\begin{equation}
 \sigma_i=\sgn(n_i),\qquad s_{ij}=\sigma_i\sigma_j.
\end{equation}
For any cycle $C=(i_1i_2\cdots i_m i_1)$,
\begin{equation}
 \prod_{r=1}^{m}s_{i_r i_{r+1}}
 =\prod_{r=1}^{m}\sigma_{i_r}\sigma_{i_{r+1}}
 =\prod_{r=1}^{m}\sigma_{i_r}^{2}=+1.
\end{equation}
Thus the baseline BC-Origin II sign graph is always balanced.

\section{Edge holonomy and vertex gauge}

Let $G=(V,E)$ be an undirected finite graph whose vertices represent shadows. Edges represent declared structural couplings.

\begin{definition}[Edge holonomy field]
A $\mathbb Z_2$ edge-holonomy field on $G$ is a symmetric assignment
\begin{equation}
 \varepsilon_{ij}=\varepsilon_{ji}\in\{-1,+1\},\qquad (ij)\in E.
\end{equation}
\end{definition}

A vertex reorientation is a choice
\begin{equation}
 g_i\in\{-1,+1\},\qquad i\in V.
\end{equation}
It acts on edge signs by
\begin{equation}
 \varepsilon_{ij}\longmapsto \varepsilon'_{ij}=g_i\varepsilon_{ij}g_j.
\end{equation}
This transformation is a finite graph analogue of a gauge change. It changes local edge representatives but not the cycle obstruction.

\begin{definition}[Cycle product]
For a cycle $C=(i_1i_2\cdots i_m i_1)$, define
\begin{equation}
 \Omega_C(\varepsilon)=\prod_{r=1}^{m}\varepsilon_{i_r i_{r+1}}.
\end{equation}
The cycle is balanced if $\Omega_C=+1$ and frustrated if $\Omega_C=-1$.
\end{definition}

\begin{proposition}[Gauge invariance of cycle products]
For every vertex reorientation $g_i\in\{-1,+1\}$,
\begin{equation}
 \Omega_C(\varepsilon')=\Omega_C(\varepsilon).
\end{equation}
\end{proposition}

\begin{proof}
Substituting $\varepsilon'_{ij}=g_i\varepsilon_{ij}g_j$ gives
\begin{equation}
 \prod_{r=1}^{m}g_{i_r}\varepsilon_{i_r i_{r+1}}g_{i_{r+1}}
 =\left(\prod_{r=1}^{m}\varepsilon_{i_r i_{r+1}}\right)
  \left(\prod_{r=1}^{m}g_{i_r}^{2}\right)
 =\Omega_C(\varepsilon),
\end{equation}
because $g_i^2=1$.
\end{proof}

\begin{theorem}[Balanced sign fields]
A $\mathbb Z_2$ edge field on a connected graph is vertex-factorizable,
\begin{equation}
 \varepsilon_{ij}=h_i h_j
\end{equation}
for some vertex signs $h_i\in\{-1,+1\}$, if and only if all cycle products are $+1$.
\end{theorem}

\begin{proof}
If $\varepsilon_{ij}=h_i h_j$, then the product around any cycle telescopes to $\prod_i h_i^2=1$.

Conversely, choose a root vertex $r$ and set $h_r=1$. For any vertex $v$, define $h_v$ as the product of edge signs along any path from $r$ to $v$. The all-cycle-positive condition makes this path-independent: two different paths form a closed cycle after reversing one path, whose product is $+1$. Therefore $h_i h_j=\varepsilon_{ij}$ for every edge.
\end{proof}

This theorem is the formal reason BC-Origin II could not generate genuine frustration using only $s_{ij}=\sigma_i\sigma_j$.

\section{Holonomy-weighted multi-shadow operator}

The complete sign entering the holonomy-weighted operator is
\begin{equation}
 \chi_{ij}=\sigma_i\varepsilon_{ij}\sigma_j.
\end{equation}
Because the baseline $\sigma_i\sigma_j$ is itself vertex-factorizable, the non-removable cycle information is carried entirely by $\varepsilon$.

The preferred BC-Origin III baseline is the gauge-clean off-diagonal operator
\begin{equation}
 (D^{\rm hol}_N)_{ii}=d_i,
\end{equation}
\begin{equation}
 (D^{\rm hol}_N)_{ij}=\eta\,\chi_{ij}K_{ij},\qquad i\ne j.
\end{equation}
The adjective ``gauge-clean'' means the following. If vertex signs are changed by $g_i$, define the diagonal orthogonal matrix
\begin{equation}
 G=\diag(g_1,\ldots,g_N).
\end{equation}
Then the transformed operator is
\begin{equation}
 D' = G D G.
\end{equation}
Therefore
\begin{equation}
 \spec(D')=\spec(D).
\end{equation}
Only gauge-invariant cycle data can affect the spectral class.

\begin{remark}[Diagonal signed-shift variant]
One can also study a convention-dependent diagonal-shift variant
\begin{equation}
 (D^{\rm diag}_N)_{ii}=d_i+\gamma\sum_{j\ne i}\chi_{ij}K_{ij}.
\end{equation}
This variant is useful as a diagnostic model, but row sums of edge signs are not gauge-invariant. The baseline spectral theorems in this paper therefore use the off-diagonal holonomy operator.
\end{remark}

\section{Minimal triangle example}

The smallest graph with possible frustration is a triangle. Take three identical shadows:
\begin{equation}
 d_1=d_2=d_3=d,
\end{equation}
and equal edge magnitude
\begin{equation}
 K_{ij}=k,\\qquad i\ne j.
\end{equation}
Absorb $\eta k$ into a single coupling scale.

\subsection{Balanced triangle}

For the balanced representative $\varepsilon_{12}=\varepsilon_{23}=\varepsilon_{13}=+1$, the operator is
\begin{equation}
 D_{\rm bal}=dI+\eta k
 \begin{pmatrix}
 0&1&1\\
 1&0&1\\
 1&1&0
 \end{pmatrix}.
\end{equation}
Its spectrum is
\begin{equation}
 \spec(D_{\rm bal})=\{d+2\eta k,\ d-\eta k,\ d-\eta k\}.
\end{equation}

\subsection{Frustrated triangle}

For a frustrated representative, choose
\begin{equation}
 \varepsilon_{12}=+1,\qquad \varepsilon_{23}=+1,\qquad \varepsilon_{13}=-1.
\end{equation}
Then
\begin{equation}
 \Omega_{123}=\varepsilon_{12}\varepsilon_{23}\varepsilon_{13}=-1.
\end{equation}
The operator is
\begin{equation}
 D_{\rm fr}=dI+\eta k
 \begin{pmatrix}
 0&1&-1\\
 1&0&1\\
 -1&1&0
 \end{pmatrix}.
\end{equation}
Its spectrum is
\begin{equation}
 \spec(D_{\rm fr})=\{d+\eta k,\ d+\eta k,\ d-2\eta k\}.
\end{equation}
Thus the frustrated cycle changes the spectral type. It is not just a relabeling of local signs.

\begin{figure}[h]
\centering
\includegraphics[width=.92\linewidth]{triangle_holonomy_spectra.png}
\caption{Balanced and frustrated triangle spectral types as the edge scale $\eta k$ is varied. The frustrated class has a different horizon margin: its lowest denominator is $d-2\eta k$ rather than $d-\eta k$.}
\end{figure}

\begin{figure}[h]
\centering
\includegraphics[width=.88\linewidth]{balanced_vs_frustrated_heatmaps.png}
\caption{Minimal triangle operators. The frustrated class cannot be converted to the balanced class by vertex reorientation because the cycle product is $-1$.}
\end{figure}

\section{Admissibility horizon and localization price}

BC-Origin I interpreted positive eigen-denominators through observable scale ratios
\begin{equation}
 S_a=\frac{\mu_{\rm eff}}{\lambda_a},\qquad \lambda_a>0.
\end{equation}
Therefore spectral splitting is not merely a drawing of lines on a graph. It changes the localization scale of observable shadows.

The structural gap between eigen-denominators can be read as a finite-dimensional localization price: hidden coupling prevents exact degeneracy, separates the spectral branches, and thereby compresses one branch while pushing another toward the horizon.

For the minimal two-shadow model
\begin{equation}
 D(t)=
 \begin{pmatrix}
 d_1(t)&\kappa\\
 \kappa&d_2(t)
 \end{pmatrix},
\end{equation}
the eigen-denominators are
\begin{equation}
 \lambda_\pm(t)=\frac{d_1(t)+d_2(t)}{2}
 \pm \sqrt{\left(\frac{d_1(t)-d_2(t)}{2}\right)^2+\kappa^2}.
\end{equation}
If $\kappa\ne0$, the branches avoid crossing. This is Landau-Zener-type avoided-crossing geometry in a finite-dimensional toy operator. It is not, by itself, a Landau-Zener transition-probability model and it is not a physical time-evolution law.

\begin{figure}[h]
\centering
\includegraphics[width=.88\linewidth]{flow_avoided_crossing_horizon.png}
\caption{Companion flow-lab intuition. A nonzero structural coupling creates avoided-crossing geometry. The lower branch may approach the admissibility horizon $\lambda=0$; this is a toy-model horizon event, not a claim of physical annihilation.}
\end{figure}

In BC-Origin III this observation is used only as interpretation. The main theorem-level content remains the graph-holonomy classification and its spectral consequences.

\section{Relation to the Flow Lab}

The software companion contains two layers.

The first layer is the BC-Origin III holonomy calculation package. It generates the triangle spectra, cycle-product diagnostics and heatmaps used in this manuscript.

The second layer is an optional BC-Origin Flow Lab. It visualizes a two-shadow parameter-flow toy model with avoided crossing and horizon crossing. This Flow Lab is a preview direction for a later BC-Origin IV article. The flow parameter should be read as a controlled deformation parameter, not automatically as physical time.

This separation prevents the present paper from mixing two claims:
\begin{enumerate}[leftmargin=2em]
\item BC-Origin III: non-factorizable edge holonomy creates genuine cycle obstruction.
\item BC-Origin Flow/IV: parameter-dependent shadow operators can show avoided crossings and horizon events.
\end{enumerate}
The first claim is the article's formal contribution. The second is a visual and conceptual companion.

\section{Controlled claims}

The paper claims:
\begin{enumerate}[leftmargin=2em]
\item The baseline BC-Origin II sign rule is balanced because it is vertex-factorizable.
\item A non-factorizable edge-holonomy field $\varepsilon_{ij}$ is the minimal $\mathbb Z_2$ graph extension required for frustrated cycle products.
\item Cycle products $\Omega_C$ are invariant under vertex reorientation.
\item A connected graph sign field is vertex-factorizable if and only if every cycle product is $+1$.
\item The holonomy-weighted off-diagonal shadow operator has gauge-invariant spectral classes under vertex sign changes.
\item The minimal triangle has different balanced and frustrated spectral types.
\item These spectral differences change toy-model admissibility-horizon margins.
\end{enumerate}

\section{Non-claims and limitations}

The paper does not claim:
\begin{enumerate}[leftmargin=2em]
\item that $\varepsilon_{ij}$ is a real physical gauge field;
\item that the model is a lattice gauge theory;
\item that it derives real spin-glass physics;
\item that spectral gaps are particle masses;
\item that horizon crossing is physical annihilation;
\item that the Flow Lab defines physical time;
\item that the model has empirical validation;
\item that it reconstructs a true hidden sector.
\end{enumerate}

The entire construction remains a finite-dimensional structural-existence model.

\section{Software appendix}

The accompanying software package has the following suggested structure:
\begin{verbatim}
bc_origin/origin_iii/
  README.md
  software/
    bc_origin_iii_holonomy_core.py
    generate_bc_origin_iii_figures.py
  figures/
    triangle_holonomy_spectra.png
    balanced_vs_frustrated_heatmaps.png
    cycle_holonomy_product.png
    flow_avoided_crossing_horizon.png

bc_origin/origin_flow/
  README.md
  web/
    index.html
\end{verbatim}

The holonomy figures are generated by
\begin{verbatim}
python software/generate_bc_origin_iii_figures.py --out figures
\end{verbatim}
from the $\texttt{bc\_origin/origin\_iii}$ folder.

The Flow Lab is a standalone browser file and can be opened directly:
\begin{verbatim}
bc_origin/origin_flow/web/index.html
\end{verbatim}

\section{Conclusion}

BC-Origin I showed that a hidden oriented residual generator can produce a toy observable shadow with scale, sign and local admissibility. BC-Origin II replaced pairwise fitted coupling by structural kernels on hidden winding-index space. BC-Origin III adds the next minimal layer: non-factorizable edge holonomy.

The main result is conceptual and finite-dimensional. If all signs are products of vertex labels, the graph is balanced and frustration is absent. If an edge-holonomy field has a negative cycle product, no vertex reorientation can remove it. This non-removable cycle data changes the spectral class of the multi-shadow operator and therefore changes the localization margins of the observable toy shadows.

The programmatic role of BC-Origin III is to open a controlled path from structural kernels to frustrated spectral geometry, while keeping the stronger dynamical interpretation for the separate Flow Lab and a later BC-Origin IV module.

\begin{thebibliography}{9}

\bibitem{BCOriginI}
A. A. Malachevsky,
\emph{Boundary Compensation Origin I: Oriented Winding Shadows, Structural Scale Selection, and Signed Spectral Localization},
BC-Origin Program preprint, 2026.

\bibitem{BCOriginII}
A. A. Malachevsky,
\emph{Boundary Compensation Origin II: Structural Coupling Kernels and Multi-Shadow Spectral Geometry},
BC-Origin Program preprint, 2026.

\bibitem{BCVIII}
A. A. Malachevsky,
\emph{Boundary Compensation VIII: Parameter Flow of Induced Gap Spectra and Sector-Weight Transport},
Boundary Compensation research programme, 2026.

\bibitem{Bollobas}
B. Bollobas,
\emph{Modern Graph Theory},
Springer, 1998.

\bibitem{Zaslavsky}
T. Zaslavsky,
Signed graphs,
\emph{Discrete Applied Mathematics} 4 (1982), 47--74.

\bibitem{Kato}
T. Kato,
\emph{Perturbation Theory for Linear Operators},
Springer, 1995.

\end{thebibliography}

\end{document}
